For a type \(\tau=(\delta,\eta)\), define its slack at \(\lambda^\dagger\) by
\[
s_{\delta,\eta}=\sum_{z=1}^{3}\lambda^\dagger_{(\eta(\delta(z)),\delta(z),z)}-\bigl(\eta(1)-\eta(0)\bigr).
\]
Using the coordinate order \(((0,0),(1,0),(0,1),(1,1))\) within every arm, direct substitution of the blocks \((2,1,2,1)\), \((0,-2,-2,0)\), and \((0,1,0,1)\) gives
\[
\begin{array}{c|rrrr}
\delta\backslash\eta&00&01&10&11\\ \hline
000&2&1&1&0\\
001&2&2&0&0\\
010&0&1&1&2\\
011&0&2&0&2\\
100&2&0&2&0\\
101&2&1&1&0\\
110&0&0&2&2\\
111&0&1&1&2
\end{array}
\]
Every slack is nonnegative, and the two gauged coordinates are zero. Thus \(\lambda^\dagger\in\mathcal H^0_{2,3}\). The zero entries in the table give exactly the twelve types displayed in the theorem statement.

Stack the ten coefficient rows indexed by \(B^\dagger\), followed by the two gauge rows, in the ambient coordinate order obtained by concatenating the three displayed arm blocks. The resulting matrix is
\[
M^\dagger=\begin{pmatrix}
0&1&0&0&0&1&0&0&0&1&0&0\\
0&1&0&0&0&1&0&0&0&0&1&0\\
0&1&0&0&0&1&0&0&0&0&0&1\\
1&0&0&0&0&0&1&0&1&0&0&0\\
1&0&0&0&0&0&1&0&0&0&1&0\\
0&1&0&0&0&0&1&0&0&0&1&0\\
0&0&0&1&1&0&0&0&1&0&0&0\\
0&0&0&1&0&1&0&0&0&1&0&0\\
0&0&1&0&0&0&1&0&1&0&0&0\\
0&0&0&1&0&0&0&1&1&0&0&0\\
0&0&0&0&1&0&0&0&0&0&0&0\\
0&0&0&0&0&0&0&0&1&0&0&0
\end{pmatrix},\qquad \det(M^\dagger)=1.
\]
Thus the first ten rows have rank \(10\), and all twelve rows have rank \(12\). Since the gauge section has dimension \(10\), these active equations isolate \(\lambda^\dagger\); hence \(\lambda^\dagger\in\mathcal V_{2,3}\).

By the gauge-transversal lemma, every \(v\in K_{2,3}\) has the form \(v_{(y,d,z)}=a_z\). Consequently
\[
I_z(\lambda+v)=\lambda_{(1,1,z)}+a_z-\lambda_{(0,1,z)}-a_z=I_z(\lambda).
\]
For a comparator indexed by \((a,s)\), where \(s=(s_0,s_1)\), treatment swapping and sign reversal give
\[
I_z(\bar w(a,s))=-w(a,s)_{1,0,z}+w(a,s)_{0,0,z}.
\]
When \(n=2\), nonconstancy implies \(s_0\ne s_1\), and the defining formula yields
\[
-w(a,s)_{1,0,z}+w(a,s)_{0,0,z}=\mathbf 1\{z=s_0\}.
\]
The gauge correction is constant within arms, so \(I(\Gamma(\bar w(a,s)))=e_{s_0}\).

It remains to verify, rather than assume from the construction, that the comparator set lies in the vertex set. For \(a=1\) and \(s=(2,3)\), the gauge-fixed comparator has blocks
\[
\mu^0=\bigl((1,0,1,1),(0,0,-1,0),(0,0,0,0)\bigr).
\]
Its slack table is
\[
\begin{array}{c|rrrr}
\delta\backslash\eta&00&01&10&11\\ \hline
000&1&0&1&0\\
001&1&0&1&0\\
010&0&0&0&0\\
011&0&0&0&0\\
100&1&0&2&1\\
101&1&0&2&1\\
110&0&0&1&1\\
111&0&0&1&1
\end{array}
\]
The ten tight types
\[
\begin{split}
\{&((000),(01)),((000),(11)),((001),(01)),((001),(11)),((010),(00)),\\
&((010),(01)),((010),(10)),((011),(00)),((100),(01)),((110),(00))\}
\end{split}
\]
together with the two gauge rows have determinant \(1\). Thus \(\mu^0\) is a vertex. Every admissible triple \((a,s_0,s_1)\) at \((n,l)=(2,3)\) is an arm permutation of \((1,2,3)\). Arm permutations preserve the quotient dual system, and \(\Gamma\) selects its unique gauge representative; hence all six comparators are vertices. They are distinct: the additional gauge-invariant contrast \(J_z(\lambda)=\lambda_{(1,0,z)}-\lambda_{(0,0,z)}\) satisfies
\[
J(\Gamma(\bar w(a,s)))=-e_a,\qquad I(\Gamma(\bar w(a,s)))=e_{s_0}.
\]
Thus the representative recovers both \(a\) and \(s_0\), while \(s_1\) is the remaining arm. Therefore \(|\mathcal F^{\mathrm{BBAK}}_{2,3}|=6\). Finally,
\[
I(\lambda^\dagger)=(1-2,0-(-2),1-0)=(-1,2,1),
\]
which is not a standard basis vector. Hence \(\lambda^\dagger\notin\mathcal F^{\mathrm{BBAK}}_{2,3}\), whereas it and all six comparator representatives belong to \(\mathcal V_{2,3}\). This proves
\[
\mathcal F^{\mathrm{BBAK}}_{2,3}\subsetneq\mathcal V_{2,3},\qquad |\mathcal V_{2,3}|\ge7.
\]